\chapter{Real-Time Risk}\label{ch:pl:real-time-risk}

Each of two desks was inside its limit all afternoon. Together they were short almost twice the firm's appetite in the same index, and
nobody saw it until the \omterm{def:pl:the-platform-map:batch}{end-of-day batch}, because nothing added the desks up while they traded. Each desk's system answered the question
it was asked --- am I inside my limit? --- and the question nobody's system was asked is the one that mattered. This chapter builds the
aggregator that asks it on every fill and every market move: exposures added up the risk hierarchy incrementally, sensitivities cached and
refreshed when the market moves, limits checked in the loop, and a what-if answer before a trade is sent.

\section{What real-time means for risk}

\begin{definition}[Real-time risk system]\label{def:pl:real-time-risk:system}
A \emph{real-time risk system}\index{real-time risk system} computes the firm's exposures and limit usage continuously from the \omterm{def:pl:position-and-pnl-services:service}{position
service}'s executions and the market's moves, fast enough that a breach is seen while the trading that causes it can still be stopped, and
serves them with their timestamps.
\end{definition}

The end-of-day risk engine of Book~6 (chapter~29) revalues every position under hundreds of scenarios once a night; the pre-trade gate of
Book~13 (chapter~22) checks one order in microseconds against limits it holds in memory. Real-time risk sits between them: seconds, not
nanoseconds or hours, and the whole firm, not one order. Its numbers are coarser than the nightly engine's (sensitivities, not full
revaluation) and its scope wider than the gate's (every desk at once). The timeliness of risk data is one of the principles of BCBS~239
(Book~6, chapter~29): a figure available the next morning describes a risk that has already been taken.

\section{Aggregation along the hierarchy}

\begin{definition}[Exposure aggregation]\label{def:pl:real-time-risk:aggregation}
\emph{Exposure aggregation}\index{exposure aggregation} adds trade-level exposures --- here delta-dollars in the index --- up the risk
hierarchy (trade, book, desk, firm) so that every node carries the sum of its descendants and can be compared with its own limit.
\end{definition}

A desk's limit and the firm's are not a partition: the firm's appetite is usually less than the sum of its desks' limits, because desks are
not expected to use their limits at the same time and in the same direction. That is exactly what they sometimes do. The chapter's
hierarchy has one firm, two desks and a book each; the firm's appetite and each desk's limit are \$50~million of index delta, which
\texttt{firm.riskctl}'s rule allows (a desk's limit never above the firm's).

\begin{omfigure}
\begin{tikzpicture}[font=\scriptsize,box/.style={draw,rounded corners=2pt,align=center,minimum width=3.4cm,minimum height=0.9cm},>=Stealth]
\node[box,fill=omThm!15] (f) at (0,0) {firm: limit \$50m\\ close $-\$90.8$m (breached)};
\node[box,fill=omDef!12] (a) at (-3,-1.5) {desk A: limit \$50m\\ close $-\$47.6$m};
\node[box,fill=omDef!12] (b) at (3,-1.5) {desk B: limit \$50m\\ close $-\$43.2$m};
\node[box,fill=gray!10] (a1) at (-3,-3) {book: index puts\\ 21\,000 units, six lines};
\node[box,fill=gray!10] (b1) at (3,-3) {book: index futures\\ 8\,800 short};
\draw[->] (a1) -- (a); \draw[->] (b1) -- (b); \draw[->] (a) -- (f); \draw[->] (b) -- (f);
\node[font=\scriptsize\itshape,align=left,anchor=west] at (5.0,-2.2) {each fill and tick adds\\ its difference upwards};
\end{tikzpicture}

\omcaption{The chapter's risk hierarchy at the close, with each node's limit and delta. Every node is inside its limit except the one no
desk's system computes.}\label{fig:pl:real-time-risk:hierarchy}
\end{omfigure}

\section{Incremental updates}

\begin{definition}[Incremental aggregation]\label{def:pl:real-time-risk:incremental}
\emph{Incremental aggregation}\index{incremental aggregation} updates aggregates by adding only the change each event causes --- a fill's
exposure, a trade's change after a market move --- to the node and its ancestors, instead of recomputing the tree from all trades.
\end{definition}

\begin{definition}[Sensitivity cache]\label{def:pl:real-time-risk:cache}
A \emph{sensitivity cache}\index{sensitivity cache} keeps each trade's sensitivities (delta, gamma, vega) computed by the pricing library at
one market snapshot, uses them to carry exposures forward as the market moves, and recomputes them when the market has moved more than a
threshold since.
\end{definition}

The chapter's aggregator (\texttt{firm.rtrisk}) keeps, for each trade, the delta-dollars it currently contributes. A fill adds the new
contribution minus the old to the trade's book and every ancestor. A market tick carries each cached delta forward with its gamma, $(\Delta +
\Gamma\,\mathrm{d}S)\,q\,S$, and applies the differences; when the spot has moved more than 0.5\% since the cache was filled, every
instrument is repriced by Book~5's library (\texttt{firm.pricing}, bump and reprice) and the differences are applied again. The aggregates are
never rebuilt; the tests check that they always equal a rebuild.

\omcode{code/firm/rtrisk/firm_rtrisk.py}{50}{76}{The sensitivity cache: Greeks from the pricing library at one spot, stale beyond a
threshold.}\label{lst:pl:real-time-risk:cache}

\omcode{code/firm/rtrisk/firm_rtrisk.py}{99}{134}{Incremental aggregation: a trade's contribution, its change added to every ancestor, and
limits checked, on every fill and every tick.}\label{lst:pl:real-time-risk:aggregator}

\begin{definition}[Risk snapshot]\label{def:pl:real-time-risk:snapshot}
A \emph{risk snapshot}\index{risk snapshot} is the aggregator's state at one time: every node's exposure, the market time it reflects, the
time and spot of the sensitivities it used and their age, and the limits breached.
\end{definition}

A snapshot that does not say how old its sensitivities are cannot be trusted; one that does can be used with judgement. At the close of the
chapter's afternoon the cache was last refreshed 65 minutes earlier (3\,890~seconds): the spot had not moved 0.5\% since, so the carried
deltas were within the threshold's accuracy.

\section{Sensitivities, full revaluation and the gap}

Carrying exposures with cached sensitivities is fast and approximate; full revaluation is exact and slow. The gap is measured by pricing
desk A's options at the close both ways for moves from $-10\%$ to $+10\%$ (\cref{fig:pl:real-time-risk:gap}): the delta-gamma P\&L from
sensitivities computed at the close against the change of the options' full values.

\begin{omfigure}
\begin{tikzpicture}
\begin{axis}[width=11.5cm,height=5.6cm,xlabel={move of the index (\%)},ylabel={error of the delta-gamma P\&L (\% of full)},xmin=-10,xmax=10,
  ymin=0,ymax=25,grid=major,grid style={gray!20},tick label style={font=\small},label style={font=\small},table/col sep=comma]
\addplot[color=omDef,very thick] table[x=move_pct,y=error_pct]{figdata/platforms/18-real-time-risk/approximation.csv};
\draw[omThm,dashed] (axis cs:-10,1) -- (axis cs:10,1) node[pos=0.5,above,font=\scriptsize]{1\%};
\draw[black,dotted] (axis cs:3.5,0) -- (axis cs:3.5,25) node[pos=0.8,right,font=\scriptsize]{$+3.5\%$};
\draw[black,dotted] (axis cs:-6.75,0) -- (axis cs:-6.75,25) node[pos=0.8,left,font=\scriptsize]{$-6.75\%$};
\end{axis}
\end{tikzpicture}

\omcaption{The error of the P\&L computed from cached delta and gamma, as a share of the full revaluation's, for desk A's 21\,000 index puts
at the close, by the size of the index move. It is zero at $\pm1\%$ (where the Greeks' bumps were), below 1\% between $-6.75\%$ and $+3.5\%$,
and grows fast on the upside, where the puts' convexity changes most. Data: \texttt{fig\_rtrisk.py}.}\label{fig:pl:real-time-risk:gap}
\end{omfigure}

Inside the range a real-time system meets in a day, the approximation is excellent: the error is exactly zero at $\pm 1\%$, because the library
computed delta and gamma from bumps of $\pm 1\%$ and the quadratic passes through them, and stays under 1\% of the full revaluation from
$-6.75\%$ to $+3.5\%$. Beyond, it fails fast --- 3.1\% at $+5\%$, 23.6\% at $+10\%$ --- which is why the cache is refreshed on moves far
smaller than those, and why stress numbers come from full revaluation (Book~6, chapter~29), not from sensitivities.

\section{Limits in the loop}

\begin{definition}[What-if check]\label{def:pl:real-time-risk:whatif}
A \emph{what-if check}\index{what-if check} computes what a proposed trade would do to the exposure of every node above its book, and which
limits it would breach, before the trade is sent, without changing the aggregator's state.
\end{definition}

\omcode{code/firm/rtrisk/firm_rtrisk.py}{144}{148}{The what-if check: the trade's contribution added to its ancestors on a copy, and the
limits it would breach.}\label{lst:pl:real-time-risk:whatif}

A \omterm{def:pl:real-time-risk:whatif}{what-if check} is the real-time system's contribution to pre-trade control: the gate of Book~13 checks an order against limits in
microseconds, the what-if tells it whether the firm-level limit has room. At the close of the chapter's afternoon, desk~B asking to sell 400
more futures would be told that its desk stays at $-\$45.2$~million, inside its limit, and the firm goes to $-\$92.8$~million, far outside.

\section{Two desks, one index}

The chapter's afternoon runs from 13:00 to 17:30 on an index starting at 5\,000, with seeded moves (20\% annual volatility and a drift of
$-1.5\%$ over the afternoon; it closes at 4\,908.5). Desk~A buys index puts --- six listed lines, three strikes and two expiries, 1\,000 units
at a time --- and desk~B sells index futures, 400 at a time, each at random moments, about every five minutes. Before each trade the desk
checks its own exposure and trades only if it stays under 90\% of its limit. Executions go through the \omterm{def:pl:position-and-pnl-services:service}{position service} of chapter~17; the
firm's aggregator adds everything up.

\begin{omfigure}
\begin{tikzpicture}
\begin{axis}[width=12cm,height=6cm,xlabel={time of day (hours)},ylabel={delta (\$ millions)},xmin=13,xmax=17.5,ymin=-100,ymax=5,
  grid=major,grid style={gray!20},tick label style={font=\small},label style={font=\small},table/col sep=comma,
  legend style={font=\footnotesize,at={(0.5,1.03)},anchor=south,legend columns=3}]
\addplot[color=omThm,very thick,const plot] table[x=hour,y=firm_m]{figdata/platforms/18-real-time-risk/exposure.csv};
\addlegendentry{firm}
\addplot[color=omDef,thick,const plot] table[x=hour,y=desk_a_m]{figdata/platforms/18-real-time-risk/exposure.csv};
\addlegendentry{desk A (puts)}
\addplot[color=omMeth,thick,dashed,const plot] table[x=hour,y=desk_b_m]{figdata/platforms/18-real-time-risk/exposure.csv};
\addlegendentry{desk B (futures)}
\draw[black,thick] (axis cs:13,-50) -- (axis cs:17.5,-50) node[pos=0.72,below,font=\scriptsize]{every limit: $-\$50$ million};
\draw[omThm,dotted,thick] (axis cs:14.397,-100) -- (axis cs:14.397,5) node[pos=0.93,right,font=\scriptsize]{firm breach 14:23:50};
\end{axis}
\end{tikzpicture}

\omcaption{Index delta of the firm and of each desk through the afternoon, in the aggregator. Neither desk goes beyond its limit --- desk~A's
largest exposure is $-\$49.3$~million, desk~B's $-\$43.6$~million --- while the firm crosses its own at 14:23:50 and ends at $-\$90.8$~million.
Data: \texttt{fig\_rtrisk.py}.}\label{fig:pl:real-time-risk:exposure}
\end{omfigure}

\begin{example}[Two desks, one index]\label{ex:pl:real-time-risk:desks}
Both desks did what their systems allowed: desk~A made 21 trades and holds 21\,000 puts, desk~B made 22 and is short 8\,800 futures; neither
exceeded its limit at any moment. The firm's limit was breached at 14:23:50, the aggregator's only alert of the afternoon; at the close the
firm is $-\$90.8$~million of delta, 1.82 times its appetite (desk~A $-\$47.6$~million, desk~B $-\$43.2$~million). Without the aggregator, the
\omterm{def:pl:the-platform-map:batch}{end-of-day batch} would have found it three hours after the breach. The cache was refreshed six times; the aggregator made 10\,570 updates.
\end{example}

\begin{remark}[Who acts on the alert]\label{rem:pl:real-time-risk:who}
A firm-level breach belongs to no desk. The aggregator's alert has to reach someone with authority over both --- the risk manager who can
block new risk-increasing orders through the gate, or reduce a desk's limit on the spot (chapter~16's \omterm{def:pl:signal-serving-and-parameter-management:store}{parameter store}, with its approvals) ---
and the limit framework has to say in advance whose risk is cut first. A real-time number that nobody may act on is only a faster report.
\end{remark}

\section{Tutorial: two desks and an aggregator}\label{tut:pl:real-time-risk:tut}

\begin{tutorial}
\textbf{Goal.} Aggregate two desks' exposures incrementally, find the firm-level breach the desks could not see, and measure the \omterm{def:pl:real-time-risk:cache}{sensitivity
cache}'s approximation. \textbf{End state:} \cref{fig:pl:real-time-risk:exposure,fig:pl:real-time-risk:gap}.
\begin{tutsteps}
\item \textbf{The afternoon}: \texttt{pl\_rtrisk.afternoon()}: the index path, the desks' trades, the aggregator's alerts.
\item \textbf{Snapshots}: \texttt{agg.snapshot(t)} at any time: exposures, cache age, breaches.
\item \textbf{What-if}: \texttt{what\_if\_at\_close(r)}.
\item \textbf{The gap}: \texttt{approximation\_error(r)} and \texttt{one\_percent\_move}.
\end{tutsteps}
\textbf{What to change next.} Lower the refresh threshold to 0.1\% and count the repricings; give desk~A calls instead of puts and see whether
the firm still breaches.
\end{tutorial}

\section{Build: the real-time aggregator}\label{bld:pl:real-time-risk:rtrisk}

\begin{build}
\bfield{Purpose} The firm's exposures and limit usage, every node of the hierarchy, updated on every fill and tick, with what-if answers for
proposed trades.
\bfield{Interface} \texttt{Hierarchy}; \texttt{SensitivityCache(md, underlying, threshold)}; \texttt{Aggregator(hierarchy, cache,
limits)} with \texttt{on\_fill}, \texttt{on\_tick}, \texttt{exposure}, \texttt{snapshot}, \texttt{what\_if}, \texttt{alerts};
\texttt{Future}; \texttt{delta\_gamma\_pnl}, \texttt{full\_pnl}.
\bfield{Rules} Aggregates are updated by differences, never rebuilt; sensitivities come from the pricing library and are refreshed beyond a
threshold; every snapshot carries the age of its sensitivities; limits are checked on every update; a what-if never changes state.
\bfield{Acceptance tests} \texttt{code/firm/rtrisk/tests/}: the hierarchy and the incremental sums against a rebuild; ticks inside the
threshold (carried with gamma) and beyond it (repriced); limits, alerts and what-if; the delta-gamma P\&L against full revaluation for a small
move.
\bfield{Stretch} Vega and rates exposures on the same hierarchy; many underlyings with a correlation-aware firm number; the aggregator fed
from the event log of chapter~24.
\end{build}

\begin{omsources}
\item One Quant Book~5, chapters 4 and 28 (Greeks, the pricing library); One Quant Book~6, chapters 21 and 29 (sensitivities, the risk
engine and its hierarchy, BCBS~239); One Quant Book~11, chapter~27, and Book~13, chapter~22 (limits and the pre-trade gate).
\end{omsources}

\section{Exercises}

\begin{exercise}[$\star$]\label{exo:pl:real-time-risk:1}
A put with delta $-0.4$ on 1\,000 units of an index at 5\,000: what is its delta in dollars, and how many such trades fit under 90\% of a
\$50~million limit?
\end{exercise}

\begin{exercise}[$\star$]\label{exo:pl:real-time-risk:2}
Why is the approximation error exactly zero at $\pm 1\%$ in \cref{fig:pl:real-time-risk:gap}?
\end{exercise}

\begin{exercise}[$\star$]\label{exo:pl:real-time-risk:3}
What does the \omterm{def:pl:real-time-risk:whatif}{what-if check} tell desk~B at the close, and why can desk~B's own system not tell it that?
\end{exercise}

\begin{exercise}[$\star\star$]\label{exo:pl:real-time-risk:4}
Why does the error grow faster on the upside than on the downside for a book of long puts?
\end{exercise}

\begin{exercise}[$\star\star$]\label{exo:pl:real-time-risk:5}
The snapshot at the close used sensitivities 65 minutes old. When is that acceptable, and what would make it not?
\end{exercise}

\begin{exercise}[$\star\star$]\label{exo:pl:real-time-risk:6}
Why is \omterm{def:pl:real-time-risk:incremental}{incremental aggregation} checked against a rebuild in the tests, and what defect would the check catch?
\end{exercise}

\begin{exercise}[$\star\star\star$]\label{exo:pl:real-time-risk:7}
\emph{Coding.} Run \texttt{afternoon} with a refresh threshold of 0.1\% and of 2\%. How many repricings does each make, and does the firm's
exposure at the close change?
\end{exercise}

\begin{exercise}[$\star\star\star$]\label{exo:pl:real-time-risk:8}
\emph{Find the flaw.} ``Every desk has a limit and a pre-trade check, so the firm cannot exceed the sum of the desks' limits, and we have set
that sum below the firm's appetite.''
\end{exercise}

\section{Problem: Two Desks, One Index}

\begin{problem}[{Weekend problem --- a breach nobody's system could see}]\label{pb:pl:real-time-risk:1}
The chapter's afternoon, desks and aggregator.

\medskip\noindent\textbf{Part I --- The system.}
\begin{enumerate}
\item What does real-time mean here, compared with the nightly engine and the pre-trade gate?
\item How are exposures aggregated, and why are desk limits not a partition of the firm's?
\item What does an incremental update add, on a fill and on a tick?
\item When is the cache refreshed, and by what?
\item What does a snapshot carry?
\end{enumerate}

\medskip\noindent\textbf{Part II --- The afternoon.}
\begin{enumerate}[resume]
\item What did each desk trade, and what did it check?
\item What were the desks' largest exposures?
\item When was the firm's limit breached?
\item What was the firm's exposure at the close, as a multiple of its appetite?
\item What does the \omterm{def:pl:real-time-risk:whatif}{what-if check} say at the close?
\end{enumerate}

\medskip\noindent\textbf{Part III --- The approximation.}
\begin{enumerate}[resume]
\item How is the cache's P\&L compared with full revaluation?
\item Between which moves is the error under 1\%?
\item What is it at $+5\%$ and at $+10\%$?
\item Why is it zero at $\pm1\%$?
\item What follows for stress numbers?
\end{enumerate}

\medskip\noindent\textbf{Part IV --- The verdict.}
\begin{enumerate}[resume]
\item State the \emph{named result}: the time at which the firm-level limit was breached while each desk stayed inside its own, the exposure at
the close, and the size of the market move beyond which the cache's P\&L error exceeds one percent of full revaluation.
\item Who should receive the alert, and what should they be able to do?
\item How would you set the refresh threshold?
\item What else would you aggregate in real time?
\item In one sentence: why must risk be added up while it is taken?
\end{enumerate}
\end{problem}

\section{Interview questions}

\begin{interviewq}[$\star$ \iqroles{developer, researcher}]\label{iq:pl:real-time-risk:1}
Two desks are each inside their limits. How can the firm be outside its own?
\end{interviewq}

\begin{interviewq}[$\star\star$ \iqroles{developer}]\label{iq:pl:real-time-risk:2}
How would you keep firm-wide exposures current as thousands of fills and ticks arrive per second?
\end{interviewq}

\begin{interviewq}[$\star\star$ \iqroles{researcher}]\label{iq:pl:real-time-risk:3}
When is a delta-gamma approximation of P\&L good enough, and when must you revalue?
\end{interviewq}

\begin{interviewq}[$\star\star$ \iqroles{developer}]\label{iq:pl:real-time-risk:4}
What should a \omterm{def:pl:real-time-risk:snapshot}{risk snapshot} tell its reader besides the numbers?
\end{interviewq}

\begin{interviewq}[$\star\star\star$ \iqroles{developer}]\label{iq:pl:real-time-risk:5}
Design a \omterm{def:pl:real-time-risk:system}{real-time risk system} for a multi-desk trading firm.
\end{interviewq}
